% ---------------------------------------------------------------------------
% Methode (francais). Source faisant autorite; les chiffres viennent de
% chiffres_fr.tex, genere depuis les JSON de resultats.
% ---------------------------------------------------------------------------
\section{Méthode}
\label{sec:methode}

\subsection{Système, mission et invariants}
\label{sec:systeme}

Le convertisseur est un abaisseur (Buck) asynchrone non idéal, \emph{physiquement
d'ordre entier}, alimenté par une source fixe $V_\mathrm{in}=48$~V, avec
$L=100~\mu$H, $C=100~\mu$F, $R=10~\Omega$, $r_L=0{,}1~\Omega$, $r_C=0{,}05~\Omega$,
$R_\mathrm{on}=0{,}05~\Omega$ et une diode de roue libre de loi affine
$v_D=V_f+r_d\,i_D$, $V_f=0{,}42$~V, $r_d=0{,}020~\Omega$. La fractionnalité
appartient exclusivement au correcteur.

La mission imposée est $24\to1\to46$~V, avec transitions à $0{,}108$ et
$0{,}216$~s et fin à $0{,}324$~s, une perturbation de charge de $+0{,}5$~A entre
$0{,}054$ et $0{,}090$~s, et des conditions initiales $(i_L,v_C)=(2{,}4~\mathrm{A},
24~\mathrm{V})$ en toutes circonstances. Le correcteur est échantillonné à
$T_s=10~\mu$s avec un retard d'un échantillon ; la MLI est à $50$~kHz et le
rapport cyclique est borné à $[0{,}02;0{,}98]$. L'anticipation nominale
\begin{equation}
d_\mathrm{ff}=\frac{(R+r_L+r_d)\,v_\mathrm{ref}}{R\,V_\mathrm{in}-(R_\mathrm{on}-r_d)\,v_\mathrm{ref}}
\label{eq:dff}
\end{equation}
est conservée sans terme en $V_f$, variables nominales figées pendant les
perturbations.

\paragraph{Coût et contraintes.}
\begin{equation}
J=\frac{1}{T}\!\int_0^T\!\frac{|r-v_o|}{\max(r,1)}\,dt
+0{,}01\,\frac{\int_0^T i_L^2\,dt}{T\cdot 20^2}
+0{,}001\,\frac{\overline{|\Delta d|}}{0{,}96},\qquad T=0{,}324~\mathrm{s},
\end{equation}
sous six contraintes $g_i\le0$ :
$g_1=I_\mathrm{max}/20-1$,
$g_2=(V_\mathrm{max}-46)/1{,}5-1$,
$g_3=|e_{24}|/0{,}05-1$,
$g_4=|e_{1}|/0{,}03-1$,
$g_5=|e_{46}|/0{,}10-1$,
$g_6=(0{,}8-V_\mathrm{min})/0{,}8$.
Les extrema sont relevés sur les sorties brutes du solveur et les erreurs
$e_\bullet$ sont des moyennes sur les fenêtres stationnaires
$[0{,}0936;0{,}1044]$, $[0{,}198;0{,}2124]$ et $[0{,}306;0{,}3204]$~s. La
\emph{marge} d'un candidat est $m=-\max_i g_i$, et sur un ensemble de cas
$m=\min_k(-\max_i g_{i,k})$ : un minimum de maxima, donc une fonction
non différentiable.

\subsection{Évaluateur commuté}
\label{sec:evaluateur}

\paragraph{Pourquoi un modèle commuté.}
L'espace de travail du modèle Simulink de référence indique que le palier de
$1$~V se situe à $3~\%$ de la frontière entre conduction continue et
discontinue, et que le saut $24\to1$~V impose environ $3{,}2$~ms de conduction
discontinue. Un évaluateur moyenné en conduction continue serait donc faux
précisément là où s'applique $g_4$. Nous simulons le convertisseur commuté
avec trois états de conduction : interrupteur passant, diode passante, et
diode bloquée ($i_L=0$), ce dernier produisant la conduction discontinue sans
modèle séparé.

\paragraph{Fronts de MLI.}
Chaque pas d'intégration (Runge--Kutta d'ordre~4) est découpé exactement sur
les fronts de la porteuse. Sans ce découpage, un pas de $1~\mu$s sur une
période de $20~\mu$s limiterait la résolution du rapport cyclique à $5~\%$,
ce qui rendrait les bornes $[0{,}02;0{,}98]$ fictives. L'écart relatif sur $J$
entre les pas de recherche ($1~\mu$s) et de confirmation ($0{,}2~\mu$s) vaut
\ChiffreEcartPas{} sur une loi de commande de référence.

\paragraph{Réalisation fractionnaire.}
L'opérateur $s^\gamma$, $\gamma\in[0;1]$, est approché par Oustaloup de
demi-ordre~3 (7 cellules) sur $[1;10^5]$~rad/s, chaque cellule
$(s+z)/(s+p)$ étant discrétisée exactement par bloqueur d'ordre zéro.
L'intégrale est réalisée comme $s^{-\lambda}=s^{-1}\cdot s^{1-\lambda}$ : un
intégrateur exact assure l'erreur statique nulle, et la cellule d'Oustaloup
porte la seule partie fractionnaire. La dérivée utilise la même réalisation,
y compris pour $\mu=1$, où la formule d'Oustaloup donne
$\omega_h(s+\omega_b)/(s+\omega_h)$, une dérivée à bande limitée.

\subsection{Synthèse FO-LQI et projection vers FO-PIDF}
\label{sec:projection}

Le modèle petit-signal est normalisé par $I=5$~A, $V=24$~V, $d=0{,}48$ et
$T_0=1$~ms, puis augmenté d'un état de mémoire. La commande LQR
($R=1$) donne $K=[K_i^x,K_v^x,K_z]$. Pour le plant linéaire nominal, sans
perturbation et hors saturation, le retour d'état se réduit algébriquement à
\begin{equation}
C_\mathrm{FOLQI}(s)=K_v^x+\frac{a}{R}
+C\,(a-r_CK_v^x)\,\frac{N s}{s+N}
+\frac{k_i}{s^\alpha},\qquad
a=K_i^x\frac{V}{I},\quad N=\frac{1}{r_CC},\quad k_i=-\frac{K_z}{T_0^{\alpha}},
\label{eq:projection}
\end{equation}
d'où l'ancrage $k_p=K_v^x+a/R$, $k_d=C(a-r_CK_v^x)$,
$\omega_f=\mathrm{clip}(N;20;2\cdot10^4)$, $b=1$, $c=0$. Le facteur
$T_0^{\alpha}$ est la voie par laquelle l'ordre parent $\alpha$ agit sur la
synthèse : l'état de mémoire est une intégrale d'ordre $\alpha$ en temps
normalisé. L'ordre intégral de la projection est ensuite remplacé par
$\lambda$, et l'ordre dérivé est $\mu$.

Dans le cas entier ($\alpha=\lambda=\mu=1$, pôle exact), la
relation~\eqref{eq:projection} reproduit la proposition~1 de l'étude
antérieure avec une erreur relative de \ChiffreErreurPropUn{} sur
$[1;10^6]$~rad/s. \textbf{C'est le seul cas exact.} La projection du pôle
$N=2\cdot10^5$~rad/s dans la bande $[20;2\cdot10^4]$~rad/s introduit à elle
seule une erreur relative de \ChiffreErreurProjPole{} en haut de bande.
L'identité FO-LQI $\to$ FO-PIDF n'est donc ni générale ni exacte ; nous la
traitons comme un ancrage, jamais comme une équivalence ni comme une preuve
d'optimalité fractionnaire.

\subsection{Audit de saturation des bornes}
\label{sec:audit}

Une campagne antérieure a consommé plus de 1800 évaluations sans jamais
atteindre la faisabilité robuste, et son diagnostic mettait en cause
l'écrêtage de $\log_{10}k_p$ sur la borne $-4$, atteinte dans $79{,}0~\%$ des
candidats. Nous traitons donc chaque borne comme un objet de mesure : pour
chaque composante de $\theta$, nous rapportons la fraction des projections
qui touchent chaque borne, sur \ChiffreAuditN{} tirages Sobol de la boîte de
décision. Nous distinguons deux natures de saturation. Une saturation
\emph{structurelle} touche une composante constante de la synthèse : la borne
ne retire alors aucun degré de liberté, elle rend seulement visible une
identité. Une saturation \emph{pathologique} touche une composante qui varie
et vient buter sur la borne : celle-ci tronque effectivement la recherche.

\subsection{Recherche étagée et six métaheuristiques}
\label{sec:recherche}

\paragraph{Vecteur de décision.}
Le vecteur v5 porte sur la \emph{synthèse} et non sur le correcteur final :
$x=[\log q_i,\log q_v,\log q_\mathrm{mem},\alpha,\mu,\log_{10}\tau_\mathrm{aw},\lambda]$,
avec $\lambda$ en septième position et $\mu$ en cinquième.

\paragraph{Méthodes.}
PSO, GA, ABC et trois variantes de projet --- pAEABC, pIGWO et pIGWO-DLH ---
reprennent les opérateurs de l'étude antérieure. Le préfixe « p » signale que
l'identité opérateur par opérateur avec les publications d'origine n'est pas
établie. Toutes reçoivent la même boîte, la même population initiale pour une
graine donnée (population 10), le même comparateur de Deb et le même budget
compté en \emph{appels d'évaluateur} ; pIGWO-DLH consommant deux appels par
comparaison, elle effectue moins d'itérations à budget égal, ce qui est voulu.

\paragraph{Comparateur.}
Faisable précède infaisable ; entre infaisables, la plus faible violation
l'emporte ; entre faisables, nous classons par \emph{marge décroissante} et non
par $J$ croissant. Le classement par $J$ produit en effet des candidats collés
à la frontière (marge nominale médiane $0{,}0066$ dans l'étude antérieure,
corrélation marge/violation robuste $-0{,}821$), soit les pires points de
départ pour la robustesse. Le critère final est lexicographique : faisabilité
sur tous les cas de conception, puis marge minimale maximale, puis pire $J$ à
marge égale ($<10^{-4}$).

\paragraph{Étages.}
L'étage~1 cherche la faisabilité nominale sur un seul scénario, avec
\ChiffreBudgetUn{} appels ; une graine qui produit moins de
\ChiffreMinFaisables{} points faisables distincts est déclarée
\emph{stérile} et rapportée comme telle. L'étage~2 cherche la robustesse sur
le jeu de conception avec \ChiffreBudgetDeux{} appels, amorcé uniquement sur
les points faisables de l'étage~1 de la \emph{même} graine. L'étage~3 rejoue
tout candidat retenu au pas de $0{,}2~\mu$s : aucun candidat n'est déclaré
faisable sans ce rejeu.

\subsection{Protocole de validation}
\label{sec:validation}

\paragraph{Ensemble d'incertitude (choix déclaré).}
Le cadrage évoque dix-sept cas robustes et seize coins sans les définir. Nous
retenons $L$, $C$ et $R$ à $\pm20~\%$ et $r_C$ à $\pm50~\%$, soit $2^4=16$
coins plus le nominal. $r_C$ est choisi parce qu'il fixe à lui seul le pôle
projeté $\omega_f=1/(r_CC)$, donc l'erreur de projection. $V_\mathrm{in}$ reste
fixe et la mission n'est jamais modifiée.

\paragraph{Trois jeux disjoints.}
Les seize coins sont partagés entre conception et validation par tirage au
sort (graine publiée \ChiffreGraineSplit), jamais après avoir observé quels
cas échouent. Le jeu de test est formé de tirages uniformes dans
l'hypercube, graine \ChiffreGraineTest{} scellée dans le protocole avant toute
évaluation, et joué une seule fois ; un fichier sentinelle interdit de le
rejouer. Un jeu consulté pour décider redevient un jeu de conception et perd
toute valeur probante.

\paragraph{Taille du test.}
Un succès sur $n$ tirages sans échec ne certifie rien au-delà de l'intervalle
de Wilson à $95~\%$. Avec $50/50$, la borne basse vaut $92{,}9~\%$. Le cadrage
indiquait qu'« environ 300 » tirages suffisent pour une borne basse de
$99~\%$ ; le calcul exact donne $98{,}74~\%$ pour $300/300$ et exige
$381$ tirages. Nous en retenons \ChiffreNTest.

\subsection{Statistiques prédéclarées}
\label{sec:stats}

L'unité statistique est la graine : trente graines communes, méthodes
appariées. Nous appliquons le test omnibus de Friedman (avec la correction
d'Iman--Davenport), puis des tests de Wilcoxon signés appariés par paires,
corrigés par la méthode de Holm, et nous rapportons la corrélation
rang-bisériale appariée comme taille d'effet. Les proportions de réussite sont
accompagnées de leur intervalle de Wilson à $95~\%$. Toutes les graines sont
rapportées, stériles comprises ; aucun seuil n'est ajusté après observation.
